% Figure 9 — xy plot, redrawn to spec [ch4.md 4.9]. The two CURVES are the
% chapter's closed-form Floor(N,k,m) (Thm.~\ref{thm:lowerbound},
% \ref{thm:split-penalty}) at m=8: the flat (joint, 2k=4) and split
% (per-reader, k=2) zero-miss floors. They are closed form by nature, not by
% omission — Floor is exactly computable, so there is no "measured" version of
% a curve to overlay.
%
% The MEASURED datum is the hollow marker at (60, 8.99): the smallest swept
% digest budget at which the R1 sweep's oracle encoder actually observed zero
% misses, from whitepaper/figures/data/r1-floor.csv (a7_experiment.py, seed
% 20260816, 4000 trials/point). It sits 3.0 bits ABOVE the analytic floor
% because the shared codebook is a randomized cover rather than an optimal
% one. That gap is the honest content of the overlay: the bound holds as a
% bound, and a real scheme pays more than it.
\begin{figure}[H]
\centering
\begin{tikzpicture}[pd figure,
  declare function={
    split(\n)  = ln(\n*(\n-1))/ln(2) - 1 - 4.80735;
    flat(\n)   = ln(\n*(\n-1)*(\n-2)*(\n-3))/ln(2) - 10.7141;
  }]
\begin{axis}[pd axis,
  width=8.6cm,height=6.0cm,
  axis lines=left,
  xmin=8,xmax=105,ymin=0,ymax=17.5,
  xtick={20,40,60,80,100},ytick={0,5,10,15},
  xlabel={swarm size $N$},
  ylabel={bits to locate},
  clip=false,axis on top,
  ]
  \draw[pd guide] (axis cs:60,0) -- (axis cs:60,15.6);
  \node[pd axis label,anchor=south] at (axis cs:60,15.7) {$N{=}60$};

  \addplot[pd rule,line width=1.05pt,domain=8:105,samples=100,smooth] {flat(x)};
  \addplot[pd focus rule,line width=1.05pt,domain=8:105,samples=100,smooth] {split(x)};

  \node[pd datum] at (axis cs:60,{flat(60)}) {};
  \node[pd focus datum] at (axis cs:60,{split(60)}) {};
  \node[pd direct label,anchor=south east] at (axis cs:59,{flat(60)+0.4}) {12.77};
  \node[pd direct label,text=hhteal,anchor=north east] at (axis cs:59,{split(60)-0.4}) {5.98};

  % measured: smallest swept budget at which the oracle encoder observed zero
  % misses (r1-floor.csv). Hollow, to read as observation against the curves.
  \node[circle,draw=hhteal,fill=hhpaper,line width=.9pt,inner sep=2.3pt]
        at (axis cs:60,8.99) {};
  \draw[pd guide] (axis cs:60,{split(60)}) -- (axis cs:60,8.75);
  \node[pd direct label,text=hhteal,anchor=west] at (axis cs:61.5,8.99)
        {8.99 measured};

  \node[pd direct label,anchor=west] at (axis cs:87,{flat(87)}) {flat digest};
  \node[pd direct label,text=hhteal,anchor=west] at (axis cs:87,{split(87)}) {split digest};
\end{axis}
\end{tikzpicture}
\caption{Bits an operator must read to locate the critical artifact, at the
lower bound of Thm.~\ref{thm:lowerbound} with $m{=}8$, $k{=}2$. One flat
digest serving both readers jointly needs $\mathrm{Floor}(N,4,8)$ bits; two
split digests, one per reader, each need only $\mathrm{Floor}(N,2,8)$. At
$N{=}60$ the split digest costs 5.98 bits against the flat digest's 12.77 ---
a 2.13$\times$ penalty, not the 2$\times$ that halving the readers would
predict. Both curves are the closed-form Floor, which is exactly computable
and so has no measured counterpart. The hollow marker is the measurement: at
$N{=}60$ the sweep's oracle encoder first achieved zero observed misses at
$8.99$ bits, three bits \emph{above} the $5.98$-bit floor, because its shared
codebook is a randomized cover and not an optimal one. Read the gap, not the
agreement --- the floor is a lower bound that holds, not a prediction a real
digest meets. No zero-miss point is marked on the flat curve because the
sweep's budget range tops out at $10.98$ bits, below that curve's $12.77$-bit
floor: the split penalty is why the measurement could not reach it
[verified, \texttt{a7\_experiment.py}, seed 20260816; data in
\texttt{figures/data/r1-floor.csv}].}
\label{fig:readpoverty}
\end{figure}
